\documentclass[11pt,a4paper]{article}

\usepackage{fontspec}
\usepackage{polyglossia}
\setmainlanguage{russian}
\setotherlanguage{english}

\IfFontExistsTF{Times New Roman}
{\setmainfont{Times New Roman}}
{\setmainfont{Libertinus Serif}}

\usepackage[a4paper,top=1.8cm,bottom=1.8cm,left=2.0cm,right=2.0cm]{geometry}
\usepackage{xcolor}
\usepackage{titlesec}
\usepackage{enumitem}
\usepackage{hyperref}

\definecolor{hseblue}{HTML}{144F7B}
\definecolor{softblue}{HTML}{EAF4FA}

\hypersetup{colorlinks=true, linkcolor=black, urlcolor=hseblue}

\setlength{\parindent}{0pt}
\setlength{\parskip}{0.48em}
\linespread{1.04}\selectfont
\setlist[itemize]{leftmargin=1.4em, topsep=2pt, itemsep=1pt}

\titleformat{\section}
  {\large\bfseries\color{hseblue}}
  {Слайд \thesection.}
  {0.5em}
  {}

\newcommand{\whatsonslide}[1]{%
  \noindent\colorbox{softblue}{%
    \parbox{\dimexpr\linewidth-2\fboxsep\relax}{\textbf{Что на слайде:} #1}%
  }\par\vspace{0.2em}
}

\begin{document}

\begin{center}
{\Large\bfseries Текст для устного рассказа презентации\par}
\vspace{0.35em}
{\large «Исследование устойчивости моделей волатильности и ценообразования опционов в различных рыночных режимах»\par}
\vspace{0.45em}
Ковалёва Мария Сергеевна \\
НИУ ВШЭ, Факультет экономических наук, Москва, 2026
\end{center}

\vspace{0.5em}
Текст рассчитан примерно на 8--12 минут спокойного выступления. Каждый раздел соответствует одному слайду финальной презентации.

\tableofcontents
\newpage

\section{Титульный слайд}

\whatsonslide{Название работы, автор, научный руководитель, факультет, программа и год защиты.}

Здравствуйте, уважаемые члены комиссии. Меня зовут Ковалёва Мария Сергеевна. Тема моей курсовой работы --- исследование устойчивости моделей волатильности и ценообразования опционов в различных рыночных режимах.

В работе я сравниваю модельные цены опционов с реальными расчётными ценами Московской биржи. Основной вопрос не в том, может ли модель идеально восстановить рынок, а в том, где и почему её ошибка становится больше: при разных режимах волатильности, сроках до экспирации и уровнях денежности опциона.

\section{Контекст и проблема}

\whatsonslide{Три блока: рынок, модели и проблема; внизу сформулирован главный вопрос о точности моделей на российских данных.}

Логика исследования начинается с простого противоречия. С одной стороны, рыночная цена агрегирует информацию, ожидания и риски участников торгов. С другой стороны, классические модели описывают цену опциона через ограниченный набор параметров: цену базового актива, страйк, срок, ставку, волатильность и тип опциона.

В реальности предпосылки моделей выполняются неполностью: волатильность меняется, доходности имеют тяжёлые хвосты, а рынок может находиться в спокойном или стрессовом состоянии. Поэтому проблема исследования такая: насколько точны модели на реальных данных российского фондового рынка и как систематически меняется их ошибка.

\section{Гипотеза и цель исследования}

\whatsonslide{Слева дана гипотеза, справа цель; внизу короткая формулировка: исследовать, от чего зависят ошибки.}

Гипотеза состоит в том, что качество аппроксимации рыночных цен модельными оценками зависит от рыночного режима и ухудшается в стрессовых состояниях рынка.

Цель работы --- исследовать устойчивость моделей волатильности и ценообразования опционов в разных рыночных режимах на основе сравнения модельных и рыночных цен. На практике это означает: я рассчитываю несколько модельных цен для одних и тех же опционов, сравниваю их с биржевыми расчётными ценами и смотрю, какие факторы объясняют ошибку: режим волатильности, срок до экспирации, денежность и способ задания волатильности.

\section{Методология исследования}

\whatsonslide{Pipeline: данные MOEX, оценка волатильности, модельное ценообразование, ошибка модели, анализ устойчивости; отдельно вынесены метрики и статистические проверки.}

На этом слайде показан общий pipeline. Сначала используются реальные данные Московской биржи: опционы серии MM на фьючерс MXI, то есть на фьючерс на Индекс МосБиржи. Период наблюдений --- с 3 января 2024 года по 19 мая 2026 года.

Далее для базового фьючерса оценивается волатильность тремя способами: историческая за 21 торговый день, историческая за 63 торговых дня и прогнозная GARCH. Затем эти значения подставляются в модели ценообразования: американскую биномиальную модель CRR и европейскую Black-76.

Ошибка определяется как модельная цена минус рыночная цена. Если она отрицательная, модель недооценивает рынок. Качество сравнивается по MAE, RMSE, средней ошибке и дополнительной метрике TVNAE, которая нормирует ошибку на временную стоимость опциона. Устойчивость результатов проверяется через статистические тесты, OLS-регрессию и подвыборки.

\section{Модели ценообразования и входные параметры}

\whatsonslide{Слева CRR как биномиальное дерево для американских опционов, справа Black-76 как аналитическая европейская формула; внизу показано, что волатильность проходит через модель и даёт цену.}

Основная модель в работе --- биномиальная модель Кокса--Росса--Рубинштейна, или CRR. Она строит дерево возможных цен базового фьючерса и на каждом шаге позволяет сравнить стоимость продолжения со стоимостью немедленного исполнения. Поэтому она подходит для американских опционов, где есть право досрочного исполнения.

Black-76 используется как европейский benchmark для опционов на фьючерсы. Она удобна аналитически, но не учитывает досрочное исполнение. Входные параметры у моделей похожи: цена фьючерса, страйк, срок до экспирации, ставка, волатильность и тип опциона. Ключевой варьируемый параметр --- волатильность: разная sigma даёт разные модельные цены и разные ошибки.

\section{Общее сравнение моделей}

\whatsonslide{Таблица MAE, RMSE и средней ошибки для трёх американских спецификаций; справа основные интерпретации.}

На общей сопоставимой выборке из 18 058 наблюдений лучшее значение MAE показывает CRR с GARCH-волатильностью: 107,0 пункта. Почти такой же результат у CRR с 63-дневной исторической волатильностью: 107,1 пункта. По RMSE немного лучше 63-дневная волатильность: 140,6 против 145,8 у GARCH. 21-дневная историческая волатильность хуже по обеим метрикам.

Важный вывод: у всех моделей средняя ошибка отрицательная. Это значит, что все три американские спецификации систематически недооценивают биржевые расчётные цены. Также видно, что короткие опционы описываются почти без ошибки, а основная сложность возникает на длинных контрактах: их трудно описать одной плоской оценкой волатильности.

\section{Ошибки по рыночным режимам}

\whatsonslide{Таблица MAE по режимам low, mid и high vol; ниже матрица средних разниц между моделями.}

Здесь показано, как ошибка меняется в зависимости от режима волатильности. В низком и среднем режимах лучшей оказывается GARCH-спецификация: MAE равен 77,5 и 118,2 пункта. Это логично: GARCH лучше улавливает текущий уровень условной волатильности.

Но в режиме высокой волатильности GARCH ухудшается до 141,0 пункта, а исторические оценки оказываются лучше: 122,8 у HV\_63d и 117,8 у HV\_21d. Это уточняет гипотезу: режим действительно связан с ошибкой, но эффект не сводится только к слову «стресс». В высоковолатильном режиме меняется состав наблюдений, в частности становится важен срок до экспирации и денежность.

\section{TVNAE}

\whatsonslide{Таблица среднего и медианного TVNAE; справа выводы о временной стоимости и коротких сроках до экспирации.}

Стандартные метрики считают ошибку по полной цене опциона. Но в выборке много длинных опционов глубоко в деньгах, где большая часть цены --- это внутренняя стоимость. Эта часть почти механически определяется разницей между фьючерсом и страйком.

Поэтому в работе вводится TVNAE --- нормированная абсолютная ошибка по временной стоимости. Временная стоимость --- это та часть цены, которую модель действительно должна объяснять через волатильность, срок и ставку.

По TVNAE лучшей становится CRR + GARCH: среднее значение 0,66 против 0,70 у HV\_63d и 0,76 у HV\_21d. Особенно сильное преимущество GARCH видно на коротких DTE 0--7 дней: там ошибка по временной стоимости примерно в три раза меньше, чем у 63-дневной исторической волатильности.

\section{CRR vs Black-76}

\whatsonslide{Таблица сравнивает американскую CRR + HV\_63d и европейскую Black-76 + HV\_63d; ниже схема премии раннего исполнения.}

На этом слайде проверяется, насколько важен сам класс модели. При одинаковой 63-дневной волатильности американская CRR даёт MAE 107,1 пункта, а европейская Black-76 --- 156,2 пункта. То есть Black-76 хуже примерно на 46 процентов по MAE и имеет более сильное отрицательное смещение.

Причина в том, что исследуемые опционы имеют американский стиль исполнения. Держатель может исполнить опцион до экспирации, а Black-76 это право не учитывает. Разность между ценой CRR и Black-76 интерпретируется как премия раннего исполнения. В среднем по выборке она равна 59,28 пункта. Следовательно, для этих инструментов американская постановка является не технической деталью, а содержательно важной частью цены.

\section{Формальные проверки и устойчивость}

\whatsonslide{Четыре блока: тест средней ошибки, Diebold--Mariano, OLS на абсолютных ошибках и robustness checks.}

Формальные проверки подтверждают основные выводы. Тест средней ошибки показывает, что у всех моделей средняя ошибка меньше нуля и статистически значима: p-value меньше 0,001. Значит, систематическое занижение не является случайным шумом.

Тесты Диболда--Мариано сравнивают модели по дневным функциям потерь. На полной выборке GARCH и HV\_63d статистически неразличимы, а HV\_21d значимо хуже HV\_63d.

OLS-регрессия показывает, что срок до экспирации и денежность объясняют вариацию ошибок сильнее, чем сам режим волатильности. R-квадрат около 0,27. Проверки устойчивости дают тот же вывод: на коротких DTE выигрывает GARCH, на длинных GARCH и HV\_63d почти равны, а изменение ставки на 12, 16 и 20 процентов не меняет качественное ранжирование.

\section{Итоговые выводы}

\whatsonslide{Пять итоговых выводов: уточнённая гипотеза, премия раннего исполнения, роль горизонта, TVNAE и систематическая недооценка рынка.}

Главный вывод такой: гипотеза не отвергается в уточнённой форме. Режим рынка связан с ошибкой, но значительную часть эффекта объясняют срок до экспирации и денежность.

Второй вывод: для рассматриваемых инструментов нужна американская модель. Black-76 даёт MAE на 46 процентов хуже, потому что не учитывает досрочное исполнение.

Третий вывод: качество модели зависит от горизонта. На коротких и средних DTE лучше работает GARCH, на длинных контрактах различия между GARCH и 63-дневной исторической волатильностью сглаживаются.

Четвёртый вывод: TVNAE показывает, что GARCH лучше описывает именно временную стоимость опциона. И пятый: все простые модели систематически недооценивают рынок, значит, одной оценки волатильности недостаточно для полного описания реальной цены.

\section{Ограничения и направления развития}

\whatsonslide{Слева ограничения исследования, справа направления развития.}

У работы есть несколько ограничений. Период анализа составляет около двух лет, рассматривается один рынок и один класс инструментов --- американские опционы на фьючерс MXI. Рыночная цена задаётся расчётной ценой биржи, а не сделочными котировками, поэтому ликвидность и bid-ask spread учтены неполностью. Кроме того, в моделях используется одна оценка волатильности на контракт и постоянная ставка.

Из этих ограничений следуют направления развития. Во-первых, можно расширить выборку по времени и рынкам. Во-вторых, построить полноценную поверхность подразумеваемой волатильности по срокам и денежности. В-третьих, добавить ликвидностные факторы: bid-ask spread, объёмы и открытый интерес. И наконец, проверить более гибкие модели: Heston, jump-модели, regime-switching или ML-корректировку модельной ошибки.

\section{Спасибо за внимание}

\whatsonslide{Финальный слайд с благодарностью.}

Спасибо за внимание. Я готова ответить на вопросы.

Если нужно коротко повторить главный результат: на российских данных по опционам на фьючерс MXI американская модель принципиально важна, GARCH полезен на коротких горизонтах и лучше описывает временную стоимость, но на полной выборке статистически близок к устойчивой 63-дневной исторической волатильности. При этом все простые модели систематически недооценивают рынок, что показывает необходимость более гибкой волатильностной структуры.

\end{document}
